\chapter{Pregled postoje\'cih re\v{s}enja i tehnologija}
\label{ch:pregled}

% Draft: English body, Serbian headings.

This chapter provides the background the rest of the thesis relies on and the set of existing
solutions the evaluation compares against. It first contrasts rasterization with ray tracing and
introduces acceleration structures and the Vulkan ray-tracing extensions (Sections~\ref{sec:rastervsrt}
to \ref{sec:vulkanrt}), then surveys existing real-time approaches to the same lighting effects
(Section~\ref{sec:postojeca}) and the denoising methods that make ray-traced effects usable at
real-time sample counts (Section~\ref{sec:denoise-pregled}).

\section{Rasterizacija naspram trasiranja zrakova}
\label{sec:rastervsrt}

Rasterization projects each triangle onto the image plane and resolves visibility with a depth
buffer. It is a stream operation that maps well to GPU hardware and excels at primary visibility and
raw throughput, but it processes one triangle at a time with no access to the rest of the scene, so
it cannot natively answer ``what else is visible from this point.'' Ray tracing takes the opposite
approach: it casts rays into the scene, finds their nearest intersections \cite{moller1997raytri}, and estimates light
transport by sampling \cite{pharr2016pbrt}. This answers arbitrary point-to-point visibility, which is what indirect
lighting, reflections, and occlusion require, at the cost of per-ray traversal work and of noise
introduced by sampling. Both are approximations of the same rendering equation
\cite{kajiya1986rendering}, which expresses the outgoing radiance $L_o$ at a point $\mathbf{x}$ in
direction $\omega_o$ as emitted radiance plus the reflected integral over the hemisphere $\Omega$:
\begin{equation}
L_o(\mathbf{x}, \omega_o) = L_e(\mathbf{x}, \omega_o) + \int_{\Omega}
f_r(\mathbf{x}, \omega_i, \omega_o)\, L_i(\mathbf{x}, \omega_i)\, (\omega_i \cdot \mathbf{n})\,
\mathrm{d}\omega_i,
\label{eq:rendering}
\end{equation}
where $f_r$ is the bidirectional reflectance distribution function (BRDF), $L_i$ the incident
radiance, $\mathbf{n}$ the surface normal, and the cosine term $(\omega_i \cdot \mathbf{n})$ accounts
for the projected solid angle. Rasterization resolves the emitted and direct terms efficiently
through projection, but the integral, which requires the incident radiance $L_i$ from every direction
and therefore scene-wide visibility, is exactly what it cannot express and what ray tracing evaluates
by sampling. A hybrid renderer uses each where it is strong \cite{akenine2018rtr}. Figure~\ref{fig:rastervsrt} contrasts the two.

\begin{figure}[tbp]
\centering
\begin{tikzpicture}[>={Stealth[]}, scale=0.9]
  \begin{scope}
    \node[font=\small] at (1.9,3.1) {Rasterizacija};
    \fill (0,1.3) circle (1.6pt) node[left, font=\scriptsize] {oko};
    \draw[thick] (1,0.35) -- (1,2.35); \node[font=\scriptsize] at (1,0.05) {ravan slike};
    \draw[fill=black!8] (3.4,0.7) -- (4.0,2.1) -- (2.8,1.95) -- cycle;
    \draw[black!55] (0,1.3) -- (3.4,0.7);
    \draw[black!55] (0,1.3) -- (4.0,2.1);
    \draw[black!55] (0,1.3) -- (2.8,1.95);
    \node[font=\scriptsize] at (3.5,2.4) {trougao};
  \end{scope}
  \begin{scope}[xshift=6.4cm]
    \node[font=\small] at (1.9,3.1) {Trasiranje zrakova};
    \fill (0,1.3) circle (1.6pt) node[left, font=\scriptsize] {oko};
    \draw[thick] (1,0.35) -- (1,2.35); \node[font=\scriptsize] at (1,0.05) {ravan slike};
    \fill[black!8] (3.5,1.3) circle (0.5);
    \draw[->, black!70] (0,1.3) -- (2.95,1.3);
    \node[font=\scriptsize] at (2.0,1.55) {zrak};
  \end{scope}
\end{tikzpicture}
\caption{Rasterizacija (projekcija trouglova na ravan slike) naspram trasiranja zrakova (zrak kroz
piksel u scenu do najbli\v{z}eg preseka).}
\label{fig:rastervsrt}
\end{figure}

\section{Strukture ubrzanja: BVH, TLAS/BLAS}
\label{sec:accel-pregled}

Testing every ray against every triangle is $O(\text{rays} \times \text{triangles})$ and hopeless at
scene scale, so ray tracing depends on a spatial acceleration structure that culls most geometry per
ray. The standard structure is a bounding-volume hierarchy (BVH): a tree of nested bounding boxes
whose traversal descends only into boxes a ray actually enters, reducing the expected cost per ray to
roughly logarithmic in the triangle count. Build quality is governed by a cost heuristic, most
commonly the surface area heuristic (SAH), which prefers splits that minimize the probability-weighted
surface area of the child nodes and thus the expected number of node and triangle tests per ray. On
current GPUs the traversal and the ray-triangle intersection are executed by dedicated
fixed-function units, which is what makes casting a ray per pixel affordable in real time. Vulkan exposes a two-level structure. A bottom-level
acceleration structure (BLAS) stores a mesh's triangle geometry in object space; a top-level
acceleration structure (TLAS) stores instances, each a transform referencing a BLAS. The split lets
many instances share one BLAS and lets the TLAS, which changes as objects move, be rebuilt per frame
while the heavier BLAS persists. A structure can be fully rebuilt (highest traversal quality) or
refit (box extents updated for small deformations, cheaper but lower quality); dynamic scenes trade
these off each frame. Figure~\ref{fig:tlasblas} shows the two-level structure and the instance
sharing it enables.

\begin{figure}[tbp]
\centering
\begin{tikzpicture}[>={Stealth[]},
  inst/.style={draw, rounded corners, font=\small, align=center, minimum width=2.4cm, minimum height=8mm},
  blas/.style={draw, rounded corners, font=\small, align=center, minimum width=3.0cm, minimum height=9mm, fill=black!5}]
  \node[inst] (i0) {Instanca 0\\(transform)};
  \node[inst, right=7mm of i0] (i1) {Instanca 1\\(transform)};
  \node[inst, right=7mm of i1] (i2) {Instanca 2\\(transform)};
  \begin{scope}[on background layer]
    \node[draw, dashed, rounded corners, fit=(i0)(i1)(i2), inner sep=7pt,
          label=above:{\small TLAS (rebuild po ,,dirty'' frejmu)}] (tlas) {};
  \end{scope}
  \node[blas, below=16mm of i0, xshift=7mm] (bA) {BLAS: Mesh A};
  \node[blas, below=16mm of i2, xshift=-7mm] (bB) {BLAS: Mesh B};
  \draw[->] (i0) -- (bA);
  \draw[->] (i1) -- (bA);
  \draw[->] (i2) -- (bB);
\end{tikzpicture}
\caption{Dvonivo struktura ubrzanja: TLAS dr\v{z}i instance (transform + pokaziva\v{c} na BLAS), pri
\v{c}emu vi\v{s}e instanci mo\v{z}e deliti isti BLAS; BLAS dr\v{z}i geometriju mre\v{z}e u lokalnom
prostoru.}
\label{fig:tlasblas}
\end{figure}

\section{Trasiranje zrakova u Vulkan-u: RT-pipeline naspram ray query}
\label{sec:vulkanrt}

Vulkan ray tracing rests on three extensions: \texttt{VK_KHR_acceleration_structure} to build and
manage the BVH, and two ways to trace against it, \texttt{VK_KHR_ray_tracing_pipeline} and
\texttt{VK_KHR_ray_query} \cite{khronos_rayquery}. The ray-tracing pipeline defines dedicated shader
stages (ray generation, closest-hit, any-hit, miss, intersection), uses a shader binding table to map
geometry to shaders, schedules rays in hardware, and supports recursion. It is the natural structure
for a full path tracer. Inline ray queries instead expose a \texttt{RayQuery} object usable from any
existing shader stage: the shader runs the traversal loop itself, with no binding table and no
recursion. Table~\ref{tab:pipeline-vs-rayquery} summarizes the difference. This thesis uses ray
queries because the effects are bounded additions to an existing rasterizing renderer rather than a
standalone tracer, as argued in Section~\ref{sec:rendergraph}.

\begin{table}[tbp]
\caption{Ray-tracing pipeline naspram inline ray query u Vulkan-u.}
\label{tab:pipeline-vs-rayquery}
\centering
\begin{tabular}{lll}
\toprule
Aspekt & RT pipeline & Ray query \\
\midrule
Shader stages       & namenski (raygen, hit, miss) & bilo koji postoje\'ci stage \\
Shader binding table & potreban & nije potreban \\
Rekurzija           & podr\v{z}ana & nije podr\v{z}ana \\
Hardversko rasporedjivanje & da & ne (poziva\v{c} pi\v{s}e petlju) \\
Integracija         & zahteva zaseban pipeline & drop-in u postoje\'ci prolaz \\
Pogodno za          & pun path tracer & ograni\v{c}eni efekti na rasteru \\
\bottomrule
\end{tabular}
\end{table}

\section{Postoje\'ca real-time re\v{s}enja za GI i RT efekte}
\label{sec:postojeca}

The effects implemented here have a long line of real-time approximations, which form the comparison
set for Chapter~\ref{ch:analiza}. \emph{Screen-space} methods compute occlusion, reflections, and
indirect light from the depth and color buffers alone: SSAO and its horizon-based refinement
\cite{mittring2007cryengine, bavoil2008hbao}, screen-space reflections, and screen-space global
illumination. They are inexpensive and widely shipped, but being limited to on-screen data they break
on off-screen occluders and under disocclusion, which is exactly the failure the ray-traced versions
avoid. \emph{Voxel} methods such as voxel cone tracing voxelize the scene and cone-trace indirect
light \cite{crassin2011vxgi}: scene-wide, but coarse, memory-heavy, and prone to leaking.
\emph{Probe} methods such as dynamic diffuse global illumination and NVIDIA's RTXGI ray-trace
irradiance into a grid of probes and interpolate \cite{majercik2019ddgi, nvidia_rtxgi}: robust and
cheap to sample, but low-frequency and sensitive to probe placement. Production \emph{hybrid} systems,
most prominently Unreal Engine's Lumen, combine software and hardware ray tracing into a complete GI
solution \cite{epic_lumen}; it is the bar this work acknowledges rather than tries to beat. Finally,
\emph{reservoir} resampling (ReSTIR and ReSTIR GI) is the current state of the art for sampling many
lights and indirect paths \cite{bitterli2020restir, ouyang2021restirgi}; it is positioned as future
work in Chapter~\ref{ch:zakljucak}.

\section{Uklanjanje \v{s}uma (denoising) za RT}
\label{sec:denoise-pregled}

The integral in Equation~\eqref{eq:rendering} has no closed form, so it is estimated by Monte Carlo
integration: drawing $N$ directions $\omega_k$ from a probability density $p$ and averaging,
\begin{equation}
L_o \approx L_e + \frac{1}{N}\sum_{k=1}^{N}
\frac{f_r(\omega_k)\, L_i(\omega_k)\, (\omega_k \cdot \mathbf{n})}{p(\omega_k)}.
\label{eq:mc}
\end{equation}
The estimator is unbiased, but its variance falls only as $1/N$, so the error (its standard
deviation) falls as $1/\sqrt{N}$: halving the noise costs four times the rays. At the one-to-few rays
per pixel a real-time budget allows, the variance is large and the raw estimate is heavily noisy, so
denoising is a mandatory part of the pipeline rather than an optional post-effect. This thesis
implements Spatiotemporal Variance-Guided Filtering (SVGF) \cite{schied2017svgf}, which combines
temporal reprojection and accumulation with a variance-guided \`a-trous spatial filter whose blur
strength adapts to the estimated per-pixel variance, so it smooths noisy regions while preserving
converged detail. Chapter~\ref{ch:implementacija} describes this engine's variant in detail. The main
alternatives are its adaptive successor A-SVGF \cite{schied2018asvgf}, NVIDIA's production Real-Time
Denoisers (ReBLUR and ReLAX) \cite{nvidia_nrd}, and machine-learning denoisers such as Intel's Open
Image Denoise \cite{intel_oidn}. Over the denoised result, temporal anti-aliasing accumulates jittered
frames to remove aliasing and to converge the remaining noise \cite{karis2014taa}; its interaction
with the per-effect denoisers is part of the implementation in Chapter~\ref{ch:implementacija}.
